\documentclass[12pt]{article}
\usepackage[utf8]{inputenc}
\usepackage[T1]{fontenc}
\usepackage{hyperref}
\title{Wik Ngathan Language Resources for K-12 Curriculum Development}
\author{Compiled by Claude Sonnet 5.5 (claude-sonnet-5-5) for Andr\'as Kornai}
\date{}
\begin{document}
\maketitle

\section{Introduction}
Wik Ngathan (ISO 639-3: wig; Glottocode: wikn1245) is a Paman (Pama-Nyungan) language of Cape York Peninsula, Queensland. It forms a dialect group with Wik-Ngatharr/Wik-Alken. Most speakers live at Aurukun. The dialects are endangered by shift to Wik-Mungkan and English. I found no school programme for it.

\section{Primary Educational Resources}
No primer, textbook or school curriculum specific to Wik Ngathan turned up in my searches.

\section{Linguistic and Reference Materials}
\begin{itemize}
\item AIATSIS Austlang record for Wik Ngathan (Y54). It rates the documentation as a grammar sketch of under 100 pages, a word list, a small text collection and some audio-visual material: \url{https://collection.aiatsis.gov.au/node/1510875}
\end{itemize}
Search results cite Peter Sutton, \emph{Wik-Ngathan Dictionary} (Adelaide: Caitlin Press, 1995), and his 1978 appendix ``The Wik-Ngathana language'' to his PhD thesis. I found no online copy or URL for either, so neither is in the JSON.

\section{Religious and Community Resources}
None specific to Wik Ngathan found.

\section{Digital and Audio Resources}
\begin{itemize}
\item ELDP archive record, ``A pan-dialectal documentation of Wik-Ngathan and Wik-Ngatharr: Cape York Peninsula, Australia'': \url{https://eldp.access.preservica.com/dk0192/}
\end{itemize}
I did not open this record, so I cannot say which items are accessible.

\section{Curriculum Development Context}
Documentation exists, mainly Sutton's work and the ELDP project, but no teaching materials were found. Wik-Mungkan is much better documented and is the neighbouring language of wider use. Any curriculum should be developed with Aurukun speakers and community organisations.
\end{document}